\section*{Capítulo \ref{ch:g9:maps-of-the-earth} --- Mapas de toda la Tierra}

\begin{solution}{exo:g9:maps-of-the-earth:1}
(a) Una \omterm{def:g9:maps-of-the-earth:projection}{proyección cartográfica}. (b) Una \omterm{def:g9:maps-of-the-earth:contour}{curva de nivel}. (c) Una \omterm{def:g9:maps-of-the-earth:numerical-scale}{escala
numérica}. (d) La \omterm{def:g9:maps-of-the-earth:contour}{equidistancia}. (e) La \omterm{def:g9:maps-of-the-earth:remote-sensing}{teledetección}.
\end{solution}

\begin{solution}{exo:g9:maps-of-the-earth:2}
(a) $4 \times 25\,000 = \num{100000}$~cm $= \qty{1}{km}$. (b) $3 \times
50\,000 = \num{150000}$~cm $= \qty{1.5}{km}$. (c) $2.5 \times 1\,000\,000
= \num{2500000}$~cm $= \qty{25}{km}$. (d) $\qty{1}{km} = \num{100000}$~cm;
$100\,000 \div 25\,000 = \qty{4}{cm}$.
\end{solution}

\begin{solution}{exo:g9:maps-of-the-earth:3}
(a) Falso: una esfera no puede aplanarse sin deformar las longitudes, así
que todo \omterm{def:g3:continents-and-oceans:world-map}{mapamundi} deforma algo. (b) Verdadero. (c) Falso: un \omterm{def:g2:drawing-to-map:map}{mapa}
equivalente conserva las superficies, pero deforma las formas, sobre todo
cerca de sus bordes. (d) Verdadero.
\end{solution}

\begin{solution}{exo:g9:maps-of-the-earth:4}
(a) La pendiente es fuerte. (b) El terreno es suave. (c) Una \omterm{def:g1:landscapes-around-us:landform}{colina}, con
su cima dentro de la curva más pequeña. (d) Un punto del terreno solo
tiene una \omterm{def:g4:world-physical-map:altitude}{altitud}, así que no puede estar en dos curvas de \omterm{def:g4:world-physical-map:altitude}{altitudes}
distintas.
\end{solution}

\begin{solution}{exo:g9:maps-of-the-earth:5}
(a) A \ang{60}\,N, el \omterm{def:g1:landscapes-around-us:water}{lago} se dibuja 4 veces demasiado grande: parece tan
grande como $4 \times 1000 = \qty{4000}{km^2}$. En el \omterm{def:g3:continents-and-oceans:equator}{ecuador} las
superficies no se agrandan, así que el \omterm{def:g1:landscapes-around-us:water}{lago} dibujado con el mismo tamaño
tiene una superficie real de \qty{4000}{km^2}. (b) $4 \div 1.33 \approx 3$
veces más grande. (c) Las superficies, y las distancias de norte a sur,
crecen sin límite hacia los \omterm{def:g3:continents-and-oceans:poles}{polos}: los \omterm{def:g3:continents-and-oceans:poles}{polos} quedarían infinitamente lejos
del \omterm{def:g3:continents-and-oceans:equator}{ecuador}.
\end{solution}

\begin{solution}{exo:g9:maps-of-the-earth:6}
(a) \qty{50}{m}. (b) La cima está a \qty{544}{m}; la curva maestra que la
rodea es la de \qty{500}{m}. (c) En el oeste: un \omterm{def:g1:landscapes-around-us:water}{río} fluye cuesta abajo,
así que su lecho está más bajo aguas abajo, al este. (d) Más o menos entre
las curvas de 300 y de \qty{400}{m}: allí la pendiente es más fuerte, y
más suave tanto cerca del fondo del \omterm{def:g1:landscapes-around-us:landform}{valle} como cerca de la cima.
\end{solution}

\begin{solution}{exo:g9:maps-of-the-earth:7}
(a) C a unos \qty{400}{m}, D a unos \qty{530}{m}. (b) Unos
\qty{135}{m}, a unos \qty{1.3}{km} de C, en el fondo del \omterm{def:g1:landscapes-around-us:landform}{valle}. (c) Entre
0.3 y \qty{0.4}{km}, aproximadamente. (d) Las alturas son pequeñas
comparadas con las distancias: sin exagerarlas, las \omterm{def:g1:landscapes-around-us:landform}{colinas} parecerían
casi planas.
\end{solution}

\begin{solution}{exo:g9:maps-of-the-earth:8}
(a) $2 \times 25\,000 = \num{50000}$~cm $= \qty{500}{m}$. (b) $300 - 200 =
\qty{100}{m}$. (c) $100 \div 500 = 0.2$, una pendiente del \qty{20}{\%}.
(d) Más suave: \qty{4}{cm} son \qty{1000}{m}, así que la pendiente es
$100 \div 1000 = \qty{10}{\%}$, la mitad de empinada.
\end{solution}

\begin{solution}{exo:g9:maps-of-the-earth:9}
(a) $14 \times 18 = 252$ órbitas. (b) $14 \times 103 = 1442$ minutos, unas
\qty{24}{h}: un día. (c) $5.5 \div 1 = 5.5$ veces. (d) Por ejemplo: un
\omterm{def:g8:ice-on-earth:glacier}{glaciar} que retrocede, un bosque talado, una ciudad que crece, una \omterm{def:g1:landscapes-around-us:water}{costa} o
un \omterm{def:g1:landscapes-around-us:water}{río} que se desplazan.
\end{solution}

\begin{solution}{exo:g9:maps-of-the-earth:10}
(a) $20\,200 \div 300\,000 \approx 0.067$~s $= 6.7 \times 10^{-2}$~s.
(b) $300\,000 \times 10^{-3} = \qty{300}{km}$. (c) Un error de tiempo
minúsculo produce un error de distancia enorme: los relojes tienen que ser
extremadamente precisos.
\end{solution}

\begin{solution}{exo:g9:maps-of-the-earth:11}
(a) Distancia a A: $\sqrt{3.6^2 + 4.8^2} = \sqrt{12.96 + 23.04} =
\sqrt{36} = 6$; distancia a B: $\sqrt{6.4^2 + 4.8^2} = \sqrt{40.96 +
23.04} = \sqrt{64} = 8$. (b) $(3.6, -4.8)$, su simétrico respecto de la
recta AB. (c) Distancia de C a $(3.6, 4.8)$: $\sqrt{3.6^2 + 5.2^2} =
\sqrt{40} \approx 6.3$; a $(3.6, -4.8)$: $\sqrt{3.6^2 + 14.8^2} =
\sqrt{232} \approx 15.2$. El receptor está en $(3.6, 4.8)$. (d) Una
distancia da un círculo (una esfera en el espacio), dos dan dos puntos
posibles; hace falta otra distancia para elegir.
\end{solution}

\begin{solution}{exo:g9:maps-of-the-earth:12}
(a) $29 \div 2.1 \approx 14$ veces. (b) $2.1 \times 10.5 \approx
22$~millones de~\unit{km^2}. (c) Groenlandia parece cubrir unos
22~millones de \unit{km^2}, cerca de los 29~millones de África: en el \omterm{def:g2:drawing-to-map:map}{mapa}
mural de Mercator, las dos parecen casi del mismo tamaño, aunque África es
14 veces mayor. (d) Una \omterm{def:g9:maps-of-the-earth:projection}{proyección equivalente}.
\end{solution}

\begin{solution}{pb:g9:maps-of-the-earth:1}
\textbf{1.} $50\,000$~cm $= \qty{500}{m}$.

\textbf{2.} $3.6 \times 500 = \qty{1800}{m} = \qty{1.8}{km}$.

\textbf{3.} B--W: $2.5 \times 500 = \qty{1250}{m}$; W--T: $2.35 \times 500
\approx \qty{1180}{m}$.

\textbf{4.} Hacia el noreste.

\textbf{5.} \qty{50}{m}.

\textbf{6.} S--B baja: $296 - 143 = \qty{153}{m}$.

\textbf{7.} B--W: $335 - 143 = \qty{192}{m}$; W--T: $544 - 335 =
\qty{209}{m}$.

\textbf{8.} $209 \div 1180 \approx 0.18$, alrededor de un \qty{18}{\%}.

\textbf{9.} Una línea recta de S a T cruzaría el \omterm{def:g1:landscapes-around-us:water}{río} donde no hay ningún
puente; el sendero baja hasta el puente en B.

\textbf{10.} $1.8 + 1.25 + 1.18 \approx \qty{4.2}{km}$ a \qty{4}{km} por
hora: unos 63 minutos; subida de $192 + 209 = \qty{401}{m}$: unos 40
minutos más. En total, alrededor de 1~h~43~min.

\textbf{11.} Desde T baja hasta B y después sube de B a S: $296 -
143 = \qty{153}{m}$.

\textbf{12.} Longitud: $1.8 + 1.25 + 1.18 \approx \textbf{4.2 km}$;
desnivel total de subida: $192 + 209 = \textbf{401 m}$.
\end{solution}
